\begin{abcliste}
\abc
Die Zeichnung (Schritt 1) stimmt, ebenso die Gleichungen in den Schritten 3 und 4. \result{\text{Schritt 2 ist falsch.}} Senkrecht wirken drei Kräfte: die Gewichtskraft nach unten, die Normalkraft und die senkrechte Komponente $F\sin\alpha$ der Zugkraft nach oben. Das Seil hebt die Kiste ein wenig an, darum drückt die Kiste weniger stark auf den Boden, und die Normalkraft ist kleiner als die Gewichtskraft. Die richtige Gleichung lautet
\[
 F_\mathrm{N} + F\sin\alpha - m\,g = 0
\]
Der Fehler wirkt sich auf die Reibung (Schritt 3) und die Beschleunigung (Schritt 4) aus.

\abc
\begin{align}
 F_\mathrm{N} &= \FNF \\
   &= \mK\cdot\g-\FZ\cdot\sin\al \\
   &= \FN \approx \result{\FNP{0}{3}}
\end{align}

\abc
Die Reibung ist $F_\mathrm{R} = \mu\,F_\mathrm{N} = \FRP{0}{3}$, und waagrecht gilt $F\cos\alpha - F_\mathrm{R} = m\,a$:
\begin{align}
 a &= \aKF \\
   &= \frac{\FZ\cdot\cos\al-\FR}{\mK} \\
   &= \aK \approx \result{\aKP{0}{2}}
\end{align}
Die Kiste beschleunigt also mehr als doppelt so stark, wie Nora berechnet hat.
\end{abcliste}
